\documentclass[12pt,a4paper]{article}
\usepackage[margin=25mm, headheight=15pt, headsep=10pt]{geometry}
\usepackage{fontspec}
\usepackage[table]{xcolor} % Note: table option is required for rowcolors
\usepackage{titlesec}
\usepackage{tocloft}
\usepackage{fancyhdr}
\usepackage{booktabs}
\usepackage{tcolorbox}
\tcbuselibrary{skins, breakable}
\usepackage[colorlinks=true, linkcolor=emerald, urlcolor=emerald, bookmarks=true]{hyperref}

\definecolor{emerald}{RGB}{0,102,51}
\definecolor{darkred}{RGB}{139,0,0}
\definecolor{lightgray}{RGB}{245,245,245}

\usepackage[nil]{babel}
\babelprovide[import, main]{bengali}
\babelprovide[import]{arabic}
\babelfont{rm}[Renderer=HarfBuzz,Scale=1.0]{Noto Serif Bengali}
\babelfont{sf}[Renderer=HarfBuzz]{Noto Sans Bengali}
\babelfont{tt}[Renderer=HarfBuzz]{Noto Sans Bengali}
\babelfont[arabic]{rm}[Renderer=HarfBuzz,Scale=1.1]{Noto Naskh Arabic}

% Section styling
\titleformat{\section}{\Large\bfseries\color{emerald}}{\thesection.}{0.5em}{}
\titleformat{\subsection}{\large\bfseries\color{emerald!85!black}}{\thesubsection.}{0.5em}{}
\titleformat{\subsubsection}{\normalsize\bfseries\color{emerald!70!black}}{\thesubsubsection.}{0.5em}{}
\titlespacing*{\section}{0pt}{18pt}{8pt}

% Fancy Header/Footer setup
\pagestyle{fancy}
\fancyhf{} % clear all header and footer fields
\fancyhead[L]{\color{emerald}\small\textbf{\nouppercase{\leftmark}}}
\fancyfoot[L]{\scriptsize\color{black!50}{\lat{AI}}-সংকলিত, আলেম-যাচাইকৃত নয়}
\fancyfoot[C]{\color{emerald!70!black}\thepage}
\renewcommand{\headrulewidth}{0.4pt}
\renewcommand{\headrule}{\hbox to\headwidth{\color{emerald}\leaders\hrule height \headrulewidth\hfill}}
\renewcommand{\footrulewidth}{0pt}

% Custom boxes
% 1. Ayat/Hadith Box
\newtcolorbox{ayatbox}[1][]{
    enhanced, breakable, 
    colback=emerald!5, 
    colframe=emerald, 
    boxrule=0.5pt, 
    arc=3pt, 
    left=10pt, right=10pt, top=10pt, bottom=10pt,
    #1
}

% 2. Quote/Translation Box
\newtcolorbox{quotebox}[1][]{
    enhanced, breakable,
    frame hidden,
    colback=lightgray,
    borderline west={3pt}{0pt}{emerald},
    left=15pt, right=10pt, top=10pt, bottom=10pt,
    sharp corners,
    #1
}

% 3. AI-generated notice box
\newtcolorbox{warnbox}[1][]{
    enhanced, breakable,
    colback=red!4,
    colframe=darkred,
    boxrule=0.8pt,
    arc=3pt,
    left=10pt, right=10pt, top=10pt, bottom=10pt,
    #1
}

% Commands
\newcommand{\arab}[1]{\foreignlanguage{arabic}{#1}}
\newcommand{\kib}[1]{\textcolor{emerald}{\textbf{#1}}}

% Latin (English) fallback: Noto Serif Bengali has NO Latin glyphs
\newfontfamily\latfont[Renderer=HarfBuzz]{Noto Serif}
\newcommand{\lat}[1]{{\latfont #1}}

\setlength{\parskip}{5pt}
\setlength{\parindent}{0pt}

% Fix for missing bullet in Noto Serif Bengali
\renewcommand{\labelitemi}{$\bullet$}



\begin{document}
\begin{center}{\Large\bfseries\color{emerald} ঘটনা ৮: উমার (রাঃ) — পানি বহন করে নিজের অহংকার ভাঙা}\end{center}
\vspace{1em}
\section*{ঘটনা ৮: উমার (রাঃ) — পানি বহন করে নিজের অহংকার ভাঙা}

\begin{warnbox}
 \textbf{গুরুত্বপূর্ণ নোটিশ (\lat{Important} \lat{Notice}):} এই কন্টেন্টটি \lat{AI} (কৃত্রিম বুদ্ধিমত্তা) দিয়ে প্রস্তুত ও সংকলিত। \lat{AI} কঠোর সূত্র-মান নীতি মেনে শুধু নির্ভরযোগ্য উৎস থেকে তথ্য সংগ্রহ করেছে — কেবল সহীহ/হাসান হাদিস ও সহীহ সনদের আসার রাখা হয়েছে, দুর্বল সনদ বাদ দেওয়া হয়েছে। তবে এটি কোনো আলেম বা মুহাদ্দিস কর্তৃক যাচাইকৃত নয়।
\end{warnbox}


\medskip\hrule\medskip

\subsection*{১. সূত্র কার্ড}

\begin{table}[h]\centering\small
\begin{tabular}{ll}
বিষয় & বিবরণ \\
\hline
\textbf{বর্ণনাকারী} & উরওয়াহ ইবনে যুবায়ের (রহ.) থেকে; উবায়দুল্লাহ ইবনে উমার থেকে \\
\textbf{গ্রন্থ} & আল-রিসালাতুল কুশায়রিয়্যাহ, পৃষ্ঠা ৩৮৪; মুসান্নাফ ইবনে আবি শায়বাহ, হাদিস ৩৭১৮০ ও ২৪১৬ \\
\textbf{অধ্যায়} & কুশায়রীর রিসালা — তাওয়াদু' (বিনয়) অধ্যায়; মুসান্নাফ — উমারের উক্তিসমূহ \\
\textbf{সনদের মান} & \textbf{সহীহ সনদের আসার (আথার)} — সাহাবি ও তাবেয়ির বর্ণনা \\
\textbf{মূল বিষয়} & ক্ষমতার অহংকার (নখওয়া) চিনে তা ভাঙার সচেতন অনুশীলন \\
\hline
\end{tabular}
\end{table}

\medskip\hrule\medskip

\subsection*{২. পটভূমি (কে, কখন, কোথায়, কেন)}

\textbf{উমার ইবনুল খাত্তাব (রাঃ)} — ইসলামের দ্বিতীয় খলিফা (\lat{caliph}) — ছিলেন একদিকে শাসক (\lat{ruler}), অন্যদিকে সাধারণ মানুষের মতো জীবনযাপনকারী। মদীনা ও পার্শ্ববর্তী অঞ্চলে তিনি নিজেই অনেক কাজ (\lat{work}) করতেন।

ঘটনা দুটি ঘটেছে উমারের খিলাফতকালে। প্রতিনিধিদল (\lat{delegation}) যখন তাঁর কাছে আসত এবং তাঁর আদেশ মেনে চলত — তখন তাঁর মনে অহংকারের (নখওয়া) সূক্ষ্ম ভাব জেগে উঠত। উমার (রাঃ) এই ভাব চিনতেন এবং সাথে সাথে তা ভাঙার জন্য ইচ্ছাকৃতভাবে (\lat{deliberately}) সাধারণ মানুষের কাজ করতেন।

\medskip\hrule\medskip

\subsection*{৩. পূর্ণ ঘটনার বর্ণনা}

\subsubsection*{প্রথম ঘটনা — প্রতিনিধিদলের পরে পানি বহন}

উরওয়াহ ইবনে যুবায়ের (রহ.) বলেন:

\begin{quote}
\textbf{\lat{“}আমি উমার ইবনুল খাত্তাব (রাঃ)-কে কাঁধে পানির মশক বহন করতে দেখেছি। আমি বললাম: "হে আমীরুল মু'মিনীন! এটা আপনার জন্য শোভনীয় নয়।"} \\
\textbf{উমার (রাঃ) বললেন: "যখন প্রতিনিধিদল আমার কাছে আসে, সকলে আমার কথা মান্য করে, তখন আমার মনে অহংকারের ভাব (নখওয়া) জেগে ওঠে। আমি চাইলাম সেটি ভেঙে ফেলি। তাই এই মশক নিয়ে আমি একজন আনসারি মহিলার ঘরে (পানি) পৌঁছে দিলাম।"\lat{”}}
\end{quote}

\subsubsection*{দ্বিতীয় ঘটনা — কন্যাকে পানি বহনে সাহায্য}

উবায়দুল্লাহ ইবনে উমার থেকে বর্ণিত:

\begin{quote}
\textbf{\lat{“}একদিন উমার (রাঃ) বসা ছিলেন। তখন একটি মেয়ে পানির মশক বহন করে যাচ্ছিল। উমার উঠে গিয়ে মেয়েটির কাছ থেকে মশকটি নিয়ে নিজের কাঁধে বহন করে তার (মেয়েটির) বাড়ি পর্যন্ত পৌঁছে দিলেন।} \\
\textbf{সাহাবিরা জিজ্ঞেস করলেন: "আল্লাহ আপনার প্রতি রহম করুন, হে আমীরুল মু'মিনীন! এটা করতে আপনাকে কী প্ররোচিত করল?"} \\
\textbf{উমার বললেন: "আমার নফস (আত্মা) নিজের উপর গর্ব করছিল (নিজেকে নিয়ে আত্মতুষ্ট — উজব)। আমি চাইলাম সেটিকে অপমানিত/অবনত করি।"\lat{”}}
\end{quote}

\medskip\hrule\medskip

\subsection*{৪. শব্দের ব্যাখ্যা}

\begin{table}[h]\centering\small
\begin{tabular}{ll}
শব্দ & অর্থ \\
\hline
\textbf{নখওয়া (\arab{نخوة})} & ক্ষমতা, পদ বা মানুষের আনুগত্যের কারণে মনে জাগা অহংকার/গর্ব \\
\textbf{কিরবাহ (\arab{قربة})} & পানির মশক/চামড়ার পাত্র \\
\textbf{নফস (\arab{نفس})} & আত্মা/সত্তা — এখানে মানুষের প্রবৃত্তি ও আত্মপ্রেম \\
\textbf{আজাবাত / আত্মতুষ্টি} & নিজের উপর সন্তুষ্টি — উজবের শুরু \\
\textbf{উজব (\arab{عجب})} & নিজের আমল, গুণ বা অবস্থাকে নিজের কৃতিত্ব মনে করা, আল্লাহর নেয়ামত না মানা \\
\hline
\end{tabular}
\end{table}

\textbf{নখওয়া ও কিবরের সম্পর্ক:} নখওয়া হলো কিবরের প্রাথমিক রূপ — ক্ষমতার ঘোরে মানুষ নিজেকে বড় ভাবা শুরু করে। উমার (রাঃ) এই সূক্ষ্ম শুরুতেই তা চিনেছেন এবং ভেঙেছেন — আগে বাড়তে না দিয়ে।

\medskip\hrule\medskip

\subsection*{৫. সংশ্লিষ্ট হাদিস ও আয়াত}

\textbf{হাদিস (মুসলিম ২৮৬৫):} নবীজি বলেছেন: \textbf{\lat{“}জাহান্নাম বলল: আমার মধ্যে অহংকারী ও উদ্ধতরা আসবে।\lat{”}} — অহংকার জাহান্নামীদের পরিচয়।

\textbf{হাদিস (মুসলিম ২৫৮৮):} নবীজি বলেছেন: \textbf{\lat{“}যে আল্লাহর জন্য বিনয় করে, আল্লাহ তাকে উন্নত করেন।\lat{”}}

\textbf{হাদিস (তিরমিজি ২০২৯):} নবীজি বলেছেন: \textbf{\lat{“}যে ব্যক্তি আল্লাহর জন্য নিজেকে ছোট করে, আল্লাহ কিয়ামতের দিন তাকে সম্মানিত করেন।\lat{”}}

\textbf{আয়াত (সূরা আল-কাসাস ২৮:৮৩):}
\begin{quote}
\textbf{\lat{“}এটি আখেরাতের ঘর, যা আমি দেবো তাদেরকে যারা পৃথিবীতে উচ্চতা (বড়ত্ব) ও বিপর্যয় কামনা করে না।\lat{”}}
\end{quote}

\textbf{আয়াত (সূরা আন-নাহল ১৬:২৩):}
\begin{quote}
\textbf{\lat{“}নিশ্চয়ই আল্লাহ অহংকারীদের পছন্দ করেন না।\lat{”}}
\end{quote}

\medskip\hrule\medskip

\subsection*{৬. রেফারেন্স}

\begin{itemize}
\item আল-রিসালাতুল কুশায়রিয়্যাহ, পৃষ্ঠা ৩৮৪
\item মুসান্নাফ ইবনে আবি শায়বাহ, হাদিস ৩৭১৮০ ও ২৪১৬ (সহীহ সনদ)
\item সহীহ মুসলিম, হাদিস ২৮৬৫ ও ২৫৮৮
\item সুনানে তিরমিজি, হাদিস ২০২৯
\item সূরা আল-কাসাস, আয়াত ৮৩; সূরা আন-নাহল, আয়াত ২৩
\end{itemize}

\end{document}
